\documentclass{article}

\usepackage[preprint]{neurips_2025}
\usepackage{float}
\usepackage{amsmath}
\usepackage{amssymb}
\usepackage{graphicx}
\usepackage{booktabs}
\usepackage[utf8]{inputenc}
\usepackage[T1]{fontenc}
\usepackage{hyperref}
\usepackage{url}
\usepackage{amsfonts}
\usepackage{nicefrac}
\usepackage{microtype}
\usepackage{xcolor}
\usepackage{tikz}
\usepackage{algorithm}
\usepackage{algpseudocode}
\usepackage{subcaption}

\newgeometry{margin=0.2in}
\begin{document}
\vspace{-2cm}
\title{\Huge \textbf{CS433: Parallel Programming \\ Assignment-2 Report (Group 8)}}

\author{
    \LARGE \textbf{Aviral Gupta (230246)}\\
    \and
    \LARGE \textbf{Aruz Awasthi (230209)}\\
    \and
    \LARGE \textbf{Yatharth Sharma (231198)}
}

\maketitle

\section{Question 1: Bounded-Buffer Producer-Consumer}

\subsection{Problem Statement}

Implement a bounded ring buffer of size $N$ shared between $P$ producers
and a single consumer. Each producer generates $n$ integer values; the
consumer reads all $P \cdot n$ values and computes their checksum.
Three synchronisation implementations are required:
\begin{itemize}
    \item \textbf{P1} -- Lock-free FIFO semaphore built from x86-64 inline assembly
    (\texttt{lock cmpxchgl}); no OS primitives.
    \item \textbf{P2} -- POSIX \texttt{sem\_t}.
    \item \textbf{P3} -- Condition variable (\texttt{pthread\_cond\_t}).
\end{itemize}
The total number of items transferred is always $P \cdot n = 100{,}000$.
Buffer sizes $N \in \{16, 64, 256\}$ and producer counts $P \in \{1, 2, 4, 8, 16, 32\}$ are swept.

\subsection{Design}

\subsubsection{Shared Ring Buffer}

The buffer is a statically allocated array of \texttt{Slot} structures,
each aligned to a 64-byte cache line to prevent false sharing. Every slot
stores a \texttt{value} and a \texttt{full} flag (0 = empty, 1 = written).
Producers atomically claim write indices with

\[
\texttt{idx} = \texttt{fetch\_and\_add(\&write\_ptr, 1)} \bmod N,
\]

where \texttt{fetch\_and\_add} is itself a CAS retry loop built on
\texttt{lock cmpxchgl}. Because $P$ producers may claim indices
concurrently and write at different speeds, the consumer cannot assume
that slot $k$ is ready as soon as a full-semaphore unit is available.
It therefore busy-spins on the per-slot \texttt{full} flag before reading.

\subsubsection{P1 -- Lock-Free FIFO Semaphore}

Two \texttt{LFSem} semaphores (\texttt{p1\_empty} initialised to $N$,
\texttt{p1\_full} to 0) control buffer capacity and item availability.
An \texttt{LFSem} uses a Michael-Scott (M\&S) lock-free FIFO queue of waiter nodes.

\begin{itemize}
    \item \textbf{Wait}: atomically decrement \texttt{count} with
    \texttt{fetch\_and\_add(\&count, -1)}. If the old value was $> 0$,
    a resource is available and the thread continues immediately (no OS call).
    Otherwise, the thread enqueues a waiter node via the M\&S FIFO tail-insertion
    and busy-waits on a per-node \texttt{wake} flag.
    \item \textbf{Post}: atomically increment \texttt{count}. If the old value
    was $\geq 0$, no thread is blocked. Otherwise, dequeue the front waiter
    (M\&S advance-head past dummy) and write \texttt{wake = 1}.
\end{itemize}

The CAS operation uses the 32-bit \texttt{lock cmpxchgl} x86 instruction with the \texttt{setzb} flag read:

\begin{verbatim}
asm("lock cmpxchgl %4, %1\nsetzb %0"
    : "=qm"(res), "+m"(*ptr), "=a"(out)
    : "a"(old_val), "r"(new_val) :);
\end{verbatim}

\subsubsection{P2 -- POSIX Semaphore}

Two \texttt{sem\_t} objects (\texttt{p2\_empty} $= N$, \texttt{p2\_full} $= 0$)
replace the lock-free semaphores. The slot-claiming logic using
\texttt{fetch\_and\_add} is identical to P1 so that multiple producers do not
overwrite the same slot. \texttt{sem\_wait} and \texttt{sem\_post} are direct
POSIX calls; on Linux these are backed by the \texttt{futex} syscall and sleep
threads in the kernel, eliminating busy-wait.

\subsubsection{P3 -- Condition Variable}

A single \texttt{pthread\_mutex\_t} protects a shared counter \texttt{p3\_count}
(items currently in the buffer) and two condition variables:
\texttt{p3\_not\_full} (producers wait here when \texttt{p3\_count == N}) and
\texttt{p3\_not\_empty} (consumer waits here when \texttt{p3\_count == 0}).
Because the mutex is held by the producer during the write, all producers are
serialised; a monotonic write pointer \texttt{p3\_wptr} is advanced inside the
lock, removing the need for \texttt{fetch\_and\_add}.

\subsection{Compilation and Usage}

\begin{verbatim}
g++ -O3 -pthread -std=c++17 q1.cpp -o q1
./q1 <N> <P> <n>
bash run.sh   # sweeps all configs, saves results_q1.txt
\end{verbatim}

\textbf{Note}: the code uses x86-64 inline assembly and must be compiled
on a Linux/x86-64 machine. Results below were collected on macOS
(where \texttt{sem\_t} uses named semaphores); timings for P2 and P3 will
differ on Linux.

\subsection{Performance Results}

\begin{figure}[H]
\centering
\includegraphics[width=\textwidth]{q1_time.png}
\caption{Q1: Execution time vs. number of producers for $N = 16, 64, 256$.
Total items always $P \cdot n = 100{,}000$.}
\label{fig:q1_time}
\end{figure}

\begin{figure}[H]
\centering
\includegraphics[width=\textwidth]{q1_speedup.png}
\caption{Q1: Speedup $T(1)/T(P)$ vs. number of producers. Dashed line marks speedup $= 1$.}
\label{fig:q1_speedup}
\end{figure}

\subsubsection{Time Tables}

\begin{table}[H]
\centering
\caption{Q1 -- Time ($\mu$s), Buffer size $N = 16$}
\begin{tabular}{r rrr}
\toprule
$P$ & P1 (lock-free) & P2 (POSIX sem) & P3 (cond var) \\
\midrule
1  & 18665 & 21659 &  44444 \\
2  & 24923 & 26046 &  31196 \\
4  & 24292 & 45282 &  59524 \\
8  & 23687 & 57129 & 101453 \\
16 & 24691 & 61456 & 142222 \\
32 & 25090 & 60575 & 142332 \\
\bottomrule
\end{tabular}
\end{table}

\begin{table}[H]
\centering
\caption{Q1 -- Time ($\mu$s), Buffer size $N = 64$}
\begin{tabular}{r rrr}
\toprule
$P$ & P1 (lock-free) & P2 (POSIX sem) & P3 (cond var) \\
\midrule
1  & 16872 & 16923 &  43796 \\
2  & 24576 & 22791 &  36382 \\
4  & 21870 & 39001 &  74318 \\
8  & 23563 & 55559 &  99135 \\
16 & 24941 & 58293 & 141759 \\
32 & 24472 & 59915 & 136922 \\
\bottomrule
\end{tabular}
\end{table}

\begin{table}[H]
\centering
\caption{Q1 -- Time ($\mu$s), Buffer size $N = 256$}
\begin{tabular}{r rrr}
\toprule
$P$ & P1 (lock-free) & P2 (POSIX sem) & P3 (cond var) \\
\midrule
1  & 18566 & 17194 &  38592 \\
2  & 25849 & 25804 &  35950 \\
4  & 20750 & 41463 &  62401 \\
8  & 23513 & 56867 & 113395 \\
16 & 24700 & 59468 & 138919 \\
32 & 24405 & 59670 & 138864 \\
\bottomrule
\end{tabular}
\end{table}

\subsubsection{Correctness}

All 18 configurations across all three methods produce checksum
$4{,}999{,}950{,}000 = \sum_{i=0}^{99999} i$, confirming that every
produced value is consumed exactly once.

\subsection{Observations and Analysis}

\paragraph{P1 -- Lock-Free Semaphore: Flat Scaling}

P1 stays near $16{,}000$--$25{,}000\,\mu$s regardless of $P$ or $N$
(Figure~\ref{fig:q1_time}). The single consumer is the throughput
bottleneck; it can consume items no faster than it reads and prints them,
independent of how many producers feed the buffer. The lock-free path
avoids all OS calls, but CAS contention on \texttt{write\_ptr} grows
with $P$, roughly cancelling the benefit of shorter per-producer work.
Buffer size $N$ barely affects P1 because the semaphore count itself
tracks available slots without blocking; the consumer's spin on the
\texttt{full} flag is the real gate, not buffer capacity.

\paragraph{P2 -- POSIX Semaphore: Sharp Degradation with $P$}

On macOS, \texttt{sem\_t} is implemented with Darwin named semaphores,
each \texttt{sem\_wait}/\texttt{sem\_post} pair issuing a system call
with kernel-scheduler round-trips. At $P=1$ P2 is comparable to P1,
but at $P=16$ the time triples to $\sim60{,}000\,\mu$s. With more
producers there are more concurrent \texttt{sem\_post} calls fighting
over the kernel's semaphore lock, and \texttt{sem\_wait} involves
thread suspension and wake-up for every single item.
\textbf{On Linux}, \texttt{sem\_t} uses the \texttt{futex} system call,
which avoids kernel involvement when no thread is blocked (fast path) --
P2 is therefore expected to be much faster and to scale far better on
the grading machine.

\paragraph{P3 -- Condition Variable: Worst at High $P$}

P3 serialises all producers through one mutex: a producer acquiring the
lock to check \texttt{p3\_count} blocks every other producer, even those
targeting unrelated slots. With $P=32$ producers contending on a single
mutex, time reaches $\sim140{,}000\,\mu$s -- roughly three times the P1
baseline. Unlike P2, there is no fast path: every producer must acquire
the mutex whether or not the buffer is full. The condition variable's
\texttt{pthread\_cond\_signal} also incurs a context switch for the
sleeping consumer or producer, adding further overhead.

\paragraph{Buffer Size $N$ Effect}

Larger $N$ reduces the probability that a producer blocks on
\texttt{p1\_empty} / \texttt{p2\_empty} / a full \texttt{p3\_count},
lowering blocking frequency. For P2 and P3 this translates to fewer
expensive OS wait-events, but the savings are modest because the consumer
(not the buffer) is the bottleneck: even an infinitely large buffer would
not allow the single consumer to process items faster.

\paragraph{Speedup Interpretation}

The speedup $T(P=1)/T(P=k)$ with fixed total items $P \cdot n = 100{,}000$
measures whether $k$ producers can feed the single consumer faster than
one producer can. For P3 at $P=2$, the speedup is $> 1$ (two producers
overlap slightly better) before convoy effects dominate at $P \geq 4$.
P1 speedup hovers near 1 across all $P$ because its bottleneck is the
consumer, not the producers.

\newpage

\section{Question 2: Concurrent Adjacency-List Graph}

\subsection{Problem Statement}

A directed graph with $N = 1{,}000$ nodes and $M = 5{,}000$ edges is
stored as an array of $N$ sorted singly-linked adjacency lists. Threads
concurrently perform \textit{search}, \textit{insert}, and \textit{delete}
operations. Three workload files with different operation mixes are used:
File~1 (80\% insert / 10\% delete / 10\% search),
File~2 (30\% insert / 30\% delete / 40\% search), and
File~3 (10\% insert / 10\% delete / 80\% search).
Six synchronisation strategies are benchmarked:

\begin{center}
\begin{tabular}{ll}
\toprule
\textbf{Program} & \textbf{Strategy} \\
\midrule
P1 & Single global \texttt{pthread\_mutex\_t} \\
P2 & One \texttt{pthread\_mutex\_t} per adjacency list (per-row) \\
P3 & One \texttt{pthread\_mutex\_t} per node, hand-over-hand (HOH) traversal \\
P4 & One \texttt{pthread\_rwlock\_t} per node, HOH traversal \\
P5 & Per-node CAS spin-lock (x86 \texttt{lock cmpxchgl}), HOH traversal \\
P6 & Optimistic CAS: lock-free search; lock + validate at write point \\
\bottomrule
\end{tabular}
\end{center}

\subsection{Design}

\subsubsection{Common Infrastructure}

All programs initialise their adjacency lists by loading the graph file
sequentially. Programs P3--P6 insert a \emph{sentinel} node with value
\texttt{INT\_MIN} as the fixed head of each list; this eliminates
head-insertion special cases during HOH traversal.

Threads are assigned a contiguous slice of the operation array by their
thread ID and execute all assigned operations.

\subsubsection{P1 -- Global Mutex}

One \texttt{pthread\_mutex\_t} protects the entire graph. Every thread
must hold the lock for any operation on any list. This completely
serialises all work and serves as the correctness baseline.

\subsubsection{P2 -- Per-List Mutex}

An array of $N$ \texttt{pthread\_mutex\_t} objects (one per row) allows
threads on different adjacency lists to proceed in full parallel.
Within one list, operations are serialised. The critical section covers
only the traversal and modification of that single list, so the lock hold
time is proportional to the list length, which is $O(M/N)$ on average.

\subsubsection{P3 -- Per-Node HOH Mutex}

Each list node carries a \texttt{pthread\_mutex\_t}. Traversal acquires
the lock on the current node before releasing the lock on the previous
node (hand-over-hand discipline). Two threads can progress concurrently
on the same list as long as they are separated by at least one node.
The sentinel head is locked first by every thread entering a list,
which makes it a serialisation bottleneck (discussed below).

\subsubsection{P4 -- Per-Node HOH RWLock}

Like P3 but replaces \texttt{pthread\_mutex\_t} with
\texttt{pthread\_rwlock\_t}. Search acquires \texttt{rdlock} (shared);
insert and delete acquire \texttt{wrlock} (exclusive). Multiple
concurrent searches through the same node proceed without blocking each
other, which greatly benefits read-heavy workloads.

\subsubsection{P5 -- Per-Node HOH CAS Spin-Lock}

Each node carries an integer \texttt{lock} field (0 = free, 1 = held).
Acquisition uses a CAS retry loop:

\begin{verbatim}
void acquire(int *lk) { while(!CAS(0, 1, lk)); }
void release(int *lk) { *lk = 0; }
\end{verbatim}

The HOH protocol is identical to P3/P4. Avoiding OS primitives reduces
per-lock overhead at the cost of busy-waiting under contention.

\subsubsection{P6 -- Optimistic CAS (Lock-Free Search)}

\textbf{Search} traverses the list entirely without acquiring any lock.
Because searches do not modify the list, they cannot corrupt state;
the list is maintained in sorted order and nodes are never freed
(explained below), so a search always terminates correctly.

\textbf{Insert/Delete} follow a two-phase approach:
\begin{enumerate}
    \item \emph{Optimistic traverse}: walk the list lock-free to locate
    the correct position (the node just before the target value).
    \item \emph{Validate and commit}: lock \texttt{prev} and \texttt{curr}
    using the CAS spin-lock. Check that (a) \texttt{prev} is not deleted,
    (b) \texttt{prev->next == curr}, and (c) \texttt{curr} has the expected
    value. If validation fails, release locks and restart from the sentinel.
\end{enumerate}
\textbf{Delete} marks \texttt{node->deleted = 1} and physically removes
the node from the list inside the same critical section. Memory is
intentionally \emph{not freed}: the lock-free search may hold a raw
pointer to the node and freeing it would cause a use-after-free. A
production system would use hazard pointers or epoch-based reclamation.

\subsection{Compilation and Usage}

\begin{verbatim}
g++ -O3 -pthread q2.cpp -o q2
./q2 <graph_file> <ops_file> <num_ops> <num_threads>
bash run.sh   # produces results.txt
\end{verbatim}

\subsection{Performance Results}

\begin{figure}[H]
\centering
\includegraphics[width=\textwidth]{q2_time.png}
\caption{Q2: Execution time vs. threads for the three operation files. Lower is better.}
\label{fig:q2_time}
\end{figure}

\begin{figure}[H]
\centering
\includegraphics[width=\textwidth]{q2_throughput.png}
\caption{Q2: Throughput (million ops/sec) vs. threads. Higher is better.}
\label{fig:q2_tp}
\end{figure}

\begin{figure}[H]
\centering
\includegraphics[width=\textwidth]{q2_speedup.png}
\caption{Q2: Speedup $T(1)/T(k)$ vs. threads. Dashed line marks speedup $= 1$.}
\label{fig:q2_speedup}
\end{figure}

\begin{figure}[H]
\centering
\includegraphics[width=\textwidth]{q2_workload_speedup.png}
\caption{Q2: Effect of workload mix on speedup for the three best-scaling strategies (P2, P4, P6).}
\label{fig:q2_workload}
\end{figure}

\subsubsection{Time Tables -- File 1 (80\% insert / 10\% delete / 10\% search)}

\begin{table}[H]
\centering
\caption{Time ($\mu$s), File 1}
\begin{tabular}{r rrrrrr}
\toprule
Threads & P1 & P2 & P3 & P4 & P5 & P6 \\
\midrule
1  &  747 & 1163 &  822 & 2301 &  850 & 1132 \\
2  & 2941 & 2115 & 2822 & 2654 & 1371 & 1962 \\
4  & 4959 & 1218 & 4534 & 1619 &  994 & 1011 \\
8  & 6067 &  795 & 6157 &  946 &  697 &  838 \\
16 & 6804 &  747 & 6838 &  850 &  725 &  691 \\
32 & 7584 &  635 & 11426 &  971 &  918 &  582 \\
\bottomrule
\end{tabular}
\end{table}

\begin{table}[H]
\centering
\caption{Speedup $T(1)/T(k)$, File 1}
\begin{tabular}{r rrrrrr}
\toprule
Threads & P1 & P2 & P3 & P4 & P5 & P6 \\
\midrule
1  & 1.00 & 1.00 & 1.00 & 1.00 & 1.00 & 1.00 \\
2  & 0.25 & 0.55 & 0.29 & 0.87 & 0.62 & 0.58 \\
4  & 0.15 & 0.95 & 0.18 & 1.42 & 0.86 & 1.12 \\
8  & 0.12 & 1.46 & 0.13 & 2.43 & 1.22 & 1.35 \\
16 & 0.11 & 1.56 & 0.12 & 2.71 & 1.17 & 1.64 \\
32 & 0.10 & 1.83 & 0.07 & 2.37 & 0.93 & 1.95 \\
\bottomrule
\end{tabular}
\end{table}

\subsubsection{Time Tables -- File 2 (30\% insert / 30\% delete / 40\% search)}

\begin{table}[H]
\centering
\caption{Time ($\mu$s), File 2}
\begin{tabular}{r rrrrrr}
\toprule
Threads & P1 & P2 & P3 & P4 & P5 & P6 \\
\midrule
1  &  765 & 1246 &  717 & 1721 &  592 &  606 \\
2  & 1567 & 1366 & 2748 & 1753 &  883 &  818 \\
4  & 2599 &  845 & 3793 & 1048 &  685 &  455 \\
8  & 3745 &  508 & 4550 &  584 &  507 &  302 \\
16 & 4071 &  406 & 4827 &  544 &  491 &  475 \\
32 & 4632 &  498 & 4734 &  534 &  458 &  444 \\
\bottomrule
\end{tabular}
\end{table}

\begin{table}[H]
\centering
\caption{Speedup $T(1)/T(k)$, File 2}
\begin{tabular}{r rrrrrr}
\toprule
Threads & P1 & P2 & P3 & P4 & P5 & P6 \\
\midrule
1  & 1.00 & 1.00 & 1.00 & 1.00 & 1.00 & 1.00 \\
2  & 0.49 & 0.91 & 0.26 & 0.98 & 0.67 & 0.74 \\
4  & 0.29 & 1.47 & 0.19 & 1.64 & 0.86 & 1.33 \\
8  & 0.20 & 2.45 & 0.16 & 2.95 & 1.17 & 2.01 \\
16 & 0.19 & 3.07 & 0.15 & 3.16 & 1.21 & 1.28 \\
32 & 0.17 & 2.50 & 0.15 & 3.22 & 1.29 & 1.36 \\
\bottomrule
\end{tabular}
\end{table}

\subsubsection{Time Tables -- File 3 (10\% insert / 10\% delete / 80\% search)}

\begin{table}[H]
\centering
\caption{Time ($\mu$s), File 3}
\begin{tabular}{r rrrrrr}
\toprule
Threads & P1 & P2 & P3 & P4 & P5 & P6 \\
\midrule
1  &  195 &  658 &  302 & 1100 & 445 & 195 \\
2  &  831 &  977 & 2678 & 1205 & 580 & 264 \\
4  & 1460 &  650 & 3251 &  732 & 503 & 157 \\
8  & 2300 &  404 & 3706 &  433 & 389 & 116 \\
16 & 2272 &  356 & 4100 &  354 & 286 & 203 \\
32 & 2790 &  508 & 4036 &  437 & 408 & 408 \\
\bottomrule
\end{tabular}
\end{table}

\begin{table}[H]
\centering
\caption{Speedup $T(1)/T(k)$, File 3}
\begin{tabular}{r rrrrrr}
\toprule
Threads & P1 & P2 & P3 & P4 & P5 & P6 \\
\midrule
1  & 1.00 & 1.00 & 1.00 & 1.00 & 1.00 & 1.00 \\
2  & 0.23 & 0.67 & 0.11 & 0.91 & 0.77 & 0.74 \\
4  & 0.13 & 1.01 & 0.09 & 1.50 & 0.88 & 1.24 \\
8  & 0.08 & 1.63 & 0.08 & 2.54 & 1.14 & 1.68 \\
16 & 0.09 & 1.85 & 0.07 & 3.11 & 1.56 & 0.96 \\
32 & 0.07 & 1.30 & 0.07 & 2.52 & 1.09 & 0.48 \\
\bottomrule
\end{tabular}
\end{table}

\subsection{Observations and Analysis}

\paragraph{P1 (Global Mutex) -- Monotonically Worse with Threads.}

P1 is the only program that degrades at every thread count.
At $T=1$ it has the lowest overhead (no per-list or per-node bookkeeping).
But every thread must acquire the same mutex before touching any node
in any list, so the graph is effectively single-threaded. Adding more
threads increases lock-acquisition contention: each \texttt{pthread\_mutex\_lock}
involves a futex syscall when the lock is held, the OS must dequeue and
schedule the next waiter, and the cache line carrying the mutex word
bounces between cores. The throughput falls from $\sim21$M ops/s at
$T=1$ to $\sim2$M ops/s at $T=32$ in File 1 -- a $10\times$ degradation.
The speedup reaches as low as $0.07\times$ (File 3, 32 threads), meaning
32 threads are 14$\times$ slower than one thread.

\paragraph{P2 (Per-List Mutex) -- Best Overall Scalability.}

P2 achieves the best consistent speedup: up to $3.07\times$ (File 2, 16 threads)
and $1.85\times$ (File 3, 16 threads). The key insight is that with $N = 1000$
rows and at most 32 threads, most concurrent operations target different rows
and therefore proceed with no interference. Only threads assigned to the same
adjacency list are serialised, and because the list is sorted with
$O(M/N) = 5$ nodes on average, the lock hold time per operation is very short.

\textbf{Why speedup $< T$:} Amdahl's Law applies. When multiple threads target
the same row (the graph is random, so occasional collisions occur), one must wait.
Thread creation, assignment scheduling, and cache misses on the mutex words add
fixed overheads. At 32 threads these overheads dominate on File 3 (read-heavy,
very short critical sections), pulling speedup down from $3.07\times$ at
16 threads to $1.30\times$ at 32.

\paragraph{P3 (HOH Mutex) -- Sentinel Hot-Spot Causes Anti-Scaling.}

P3's speedup is consistently $< 1$ beyond two threads (as low as $0.07\times$
at 32 threads, File 3). The culprit is the sentinel head node.

In the HOH protocol, every thread entering any list must first lock the sentinel.
With $T$ threads all assigned to the same row, they form a queue at the sentinel,
re-creating the global serialisation of P1 within each row. Worse, HOH requires
holding two locks simultaneously (prev and curr), doubling lock traffic and
memory contention. The cost of \texttt{pthread\_mutex\_lock}/\texttt{unlock}
per node exceeds the cost of the actual work (a pointer comparison),
making synchronisation overhead the dominant term.

\textbf{Why P3 at $T=1$ is faster than P4/P5:}
At a single thread, P3 performs no actual locking (the mutex is always
immediately available), so the overhead is just a compare-and-swap
returning success. P4 and P5 do the same, but rwlock and CAS have
slightly higher single-threaded setup cost.

\paragraph{P4 (HOH RWLock) -- Good Scaling, Read-Dominated.}

P4 fixes P3's main weakness for search: multiple threads holding a
\texttt{rdlock} on the same node do not block each other. On File 1
(90\% writes), P4 reaches $2.71\times$ speedup at 16 threads; on
File 2 (40\% searches), it peaks at $3.22\times$ at 32 threads.

\textbf{Why P4 degrades at 32 threads (File 1):} POSIX rwlocks on Linux
are generally writer-preferring. With many concurrent writers competing
for exclusive locks, write requests block incoming read requests, creating
writer convoys. The rwlock's internal state word also causes more cache
invalidations than a simple spin-lock.

\textbf{Why P4 at $T=1$ is slower than P1/P3:} A \texttt{pthread\_rwlock\_t}
is larger and more complex than a \texttt{pthread\_mutex\_t}. The initial
state transition from free to rdlock-held involves an atomic read-modify-write;
this dominates when the critical section is a single pointer comparison.

\paragraph{P5 (HOH CAS Spin-Lock) -- Lightweight, Moderate Gains.}

P5 replaces OS-managed locks with a simple two-word CAS spin-lock.
Acquisition is a tight CPU loop on a cached cache line, avoiding context
switches. At low contention (few threads per row), P5 is competitive:
$1.22\times$ speedup at 8 threads in File 1. At high thread counts the
spin-wait loops waste CPU cycles and increase cache-coherence traffic
without doing useful work, limiting gains to $\sim1.3\times$.

Compared with P3: P5 uses the same HOH protocol but the spin-lock is
cheaper per operation, so P5 consistently outperforms P3 at every thread
count. Compared with P2: P5 still serialises all threads on the same row
through the sentinel spin, whereas P2 needs only one lock acquisition per
operation (not per node).

\paragraph{P6 (Optimistic CAS) -- Best for Read-Heavy, Degrades on Write-Heavy.}

P6 achieves the highest throughput on read-heavy workloads: $116\,\mu$s
at 8 threads in File 3 (137.9M ops/s), versus $404\,\mu$s for P2. The
lock-free search path allows every thread to traverse nodes without any
synchronisation cost, making searches theoretically linear-scalable.

\textbf{Why P6 degrades at 32 threads (Files 1, 3):} Under high write
concurrency, the optimistic validation frequently fails. A thread
optimistically traverses the list, then finds that \texttt{prev->next}
has been changed by another inserter; it must restart from the sentinel.
In the worst case, every write attempt restarts several times, resembling
livelock. This is dramatically visible in File 3 at 32 threads
($408\,\mu$s vs $116\,\mu$s at 8 threads -- a $3.5\times$ slowdown from peak).

\textbf{Why P6 at $T=1$ matches P1 (File 3):} With one thread there is no
contention so validations never fail; the lock-free traversal costs less
per node than P2--P5 (no lock operations at all for searches). File 3 is
80\% search, so the insert/delete paths are rarely used, making P6's
$195\,\mu$s match P1's $195\,\mu$s exactly.

\paragraph{Effect of Workload Mix on Speedup (Figure~\ref{fig:q2_workload}).}

For all programs, File 3 (read-heavy) shows the highest speedup for
read-allowing strategies (P4, P6) because more operations can proceed in
parallel. File 1 (write-heavy) suppresses speedup for optimistic P6
(frequent CAS retries) but not for lock-based P2/P4 (locks serialise
writes but do not retry). File 2 (balanced) sits between the two extremes
and shows the steadiest scaling for P2 and P4.

\paragraph{Summary Table.}

\begin{table}[H]
\centering
\caption{Peak speedup over all files and thread counts per strategy}
\begin{tabular}{llrr}
\toprule
Method & Key mechanism & Peak speedup & Best workload \\
\midrule
P1 & Global mutex      & $0.25\times$ & --- (always degrades) \\
P2 & Per-list mutex    & $3.07\times$ & Balanced / write-heavy \\
P3 & HOH mutex         & $0.29\times$ & --- (sentinel bottleneck) \\
P4 & HOH rwlock        & $3.22\times$ & Balanced / read-heavy \\
P5 & HOH CAS spin      & $1.56\times$ & Read-heavy, low $T$ \\
P6 & Optimistic CAS    & $2.01\times$ & Read-heavy, $T \leq 8$ \\
\bottomrule
\end{tabular}
\end{table}

\end{document}
